% =====================================================================
%  FRAGMENTOS DE CORRECCIÓN (texto listo para pegar)
%  Cada bloque indica el archivo y el pasaje que reemplaza. Ningún original
%  fue modificado. Lo que requiere un dato que solo la autora tiene está
%  marcado con \pendiente{}.
% =====================================================================

% ---------------------------------------------------------------------
% F1. Realidad aumentada: alinear el uso del término con Azuma (1997)
%     Reemplaza el pasaje de capitulos/cap4/produccion.tex, párrafo
%     «Diseño» del escáner (líneas ~315-317: «Esta superposición ...
%     constituye la realidad aumentada del sistema, sin anclaje
%     tridimensional»).
%     NOTA: el título del trabajo NO se cambia aquí; modificar el perfil
%     requiere el trámite del Cap. IX del Reglamento. Se acota el término.
% ---------------------------------------------------------------------
Esta superposición de información digital sobre la imagen en vivo de la
cámara corresponde a una forma básica de realidad aumentada: combina lo real y
lo virtual y responde a la acción del usuario, pero no cumple la tercera
característica definida por Azuma (1997), el registro en tres dimensiones,
porque no emplea seguimiento de superficies ni ARCore. En esta memoria, el
término «realidad aumentada» se emplea en ese sentido acotado
(superposición informativa sobre la cámara); el anclaje tridimensional se
propone como trabajo futuro (sección~\ref{sec:recomendaciones}).

% ---------------------------------------------------------------------
% F2. Mobile-D: precisar la fuente y la adaptación de las fases
%     Insertar al inicio de capitulos/03-area-aplicacion.tex, sección
%     «Aplicación de la Metodología Mobile-D».
%     Abrahamsson et al. (2004) estructuran el trabajo en fases y en tres
%     tipos de día: planificación, trabajo y entrega.
% ---------------------------------------------------------------------
Mobile-D (Abrahamsson et~al., 2004) organiza el desarrollo en fases y en tres
tipos de día: de planificación, de trabajo y de entrega. En este proyecto, las
fases se denominan Exploración, Inicialización, Producción, Estabilización y
Pruebas, y los días de planificación, de trabajo y de entrega se aplicaron
dentro de cada iteración de la fase de Producción. \pendiente{verificar en el
texto base de Mobile-D que la denominación y el orden de las cinco fases
coinciden con las de la fuente citada, o declarar expresamente que son una
adaptación; la versión consultada del artículo de Abrahamsson et~al. (2004)
nombra las fases \emph{set-up}, \emph{core}, \emph{core2}, \emph{stabilize} y
\emph{wrap-up}.}

% ---------------------------------------------------------------------
% F3. Arquitectura: servicios del lado del servidor omitidos en la memoria
%     Insertar en la descripción de la arquitectura (cap. 4, fase de
%     Inicialización o Producción) y actualizar el diagrama.
%     Fuente: functions/src/index.mjs, functions/src/quota.mjs,
%     src/services/firebase/firebase.ts, firestore.rules.
% ---------------------------------------------------------------------
Además de los servicios de autenticación, base de datos y almacenamiento de
Firebase, el sistema incluye dos Cloud Functions invocables
(\archivo{buscarAlimentos} y \archivo{guardarRutina}), desplegadas en la región
\archivo{us-central1} con un máximo de cinco instancias y con verificación de
Firebase App Check obligatoria. La función \archivo{buscarAlimentos} actúa como
\emph{proxy} de FoodData Central: la clave de la API se guarda como secreto del
servidor y nunca se incluye en la aplicación. Ambas funciones rechazan a los
usuarios anónimos y aplican cuotas por usuario (15 búsquedas por minuto y 120
por hora para alimentos; 20 solicitudes por minuto y 200 por hora para
rutinas); la búsqueda de alimentos añade un límite global de 800 solicitudes por
hora. Las cuotas se controlan mediante transacciones sobre documentos de
Firestore cuyo identificador es un \emph{hash}, sin almacenar \emph{tokens}.
Las reglas de seguridad de Firestore validan campos permitidos y rangos del perfil
(por ejemplo, altura de 50 a 300~cm, peso de 10 a 500~kg y edad de 14 a 90
años). La aplicación fuerza el \emph{long polling} de Firestore y no habilita
la persistencia sin conexión.

% ---------------------------------------------------------------------
% F4. Ética, privacidad y tratamiento de datos (sección nueva)
%     Insertar como sección en el capítulo 4 (Estabilización) o como anexo.
% ---------------------------------------------------------------------
\section{Consideraciones éticas y de privacidad}
\label{sec:etica}

La aplicación almacena datos personales de salud (peso, estatura, edad, género
y objetivo de entrenamiento) asociados a una cuenta de correo electrónico. Para
protegerlos, las reglas de Firestore restringen la lectura y la escritura al
propietario de cada documento, y el acceso a los servicios del servidor exige
App Check. Dado que el perfil acepta edades desde los 14 años, se requiere
\pendiente{definir y documentar la política de privacidad, el aviso de
consentimiento informado, el tratamiento de usuarios menores de edad y el
marco legal aplicable; esta auditoría no verificó la normativa boliviana de
protección de datos}. Asimismo, el sistema \textbf{no reemplaza} la supervisión
de entrenadores ni de nutricionistas, y las estimaciones de calorías e IMC son
orientativas.

% ---------------------------------------------------------------------
% F5. Licencias y atribuciones de contenido de terceros (tabla de anexo)
%     La fase de Exploración ya evalúa licencias; falta una tabla final de
%     atribuciones y autorizaciones. Completar las celdas \pendiente.
% ---------------------------------------------------------------------
\begin{table}[H]
\centering
\caption{Contenido de terceros, licencia y atribución}
\label{tab:licencias-final}
\small
\begin{tabularx}{\textwidth}{|L{3.4cm}|Y|L{3.6cm}|}
\hline
\rowcolor{encabezado} \textbf{Recurso} & \textbf{Condición de uso} & \textbf{Estado} \\
\hline
FoodData Central (USDA) & Datos del gobierno de EE.\,UU. consultados mediante
un \emph{proxy} con clave en el servidor. & \pendiente{citar los términos de uso vigentes} \\ \hline
Animaciones de ejercicios (95) & Obtenidas por la vía oficial del proveedor;
no deben usarse para entrenar modelos sin permiso escrito. &
\pendiente{adjuntar la autorización o plan vigente} \\ \hline
Fotografías del catálogo (399) & Usadas para entrenar el reconocedor. &
\pendiente{obtener la autorización del proveedor antes de publicar en Google
Play} \\ \hline
Modelo MobileNetV2 preentrenado & Referenciado en Sandler et~al. (2018). &
\pendiente{indicar la licencia de los pesos usados} \\
\hline
\end{tabularx}
\fuente{Capítulo~\ref{cap:desarrollo}, fase de exploración, y capítulo de conclusiones.}
\end{table}

% ---------------------------------------------------------------------
% F6. Alcance: separar «conteo calórico» de «calculadora de IMC»
%     Reemplaza el ítem de capitulos/01-introduccion.tex, sección Alcance
%     («El registro de alimentación se limitará solo al conteo calórico,
%     es decir una calculadora de IMC.»). El cambio del perfil aprobado
%     debe consultarse con el tutor.
% ---------------------------------------------------------------------
    \item El registro de alimentación se limitará al conteo calórico de los
    alimentos consumidos; el sistema incluye además una calculadora del índice
    de masa corporal (IMC) para el seguimiento del progreso físico.

% ---------------------------------------------------------------------
% F7. Evidencia sobre adherencia (reemplaza los 3 pasajes con Gjestvang)
%     capitulos/01-introduccion.tex líneas ~5-9, ~39-43 y ~79-83.
%     Verificado: la cifra «40-65 %» NO figura en Gjestvang et al. (2020).
%     Verificado en la introducción de Gjestvang et al. (2021): menos del
%     40 % de los socios mantiene ejercicio regular a largo plazo, y la
%     cancelación de membresía fue baja (1,8 % a los 3 meses, 5,2 % a los 6
%     y 13,4 % a los 12 meses). Ver nota de pie para las barreras.
% ---------------------------------------------------------------------
Gjestvang et~al. (2020) realizaron un estudio de seguimiento de un año con
nuevos socios de centros de fitness en Noruega para identificar los motivos y
las barreras para iniciar y sostener el ejercicio. Investigaciones del mismo
grupo señalan que menos del 40\,\% de los socios de gimnasios logra mantener un
ejercicio regular a largo plazo, aunque pocos cancelan su membresía (13,4\,\% a
los 12 meses) (Gjestvang et~al., 2021). Estos datos describen un problema de
adherencia, no de abandono formal de la membresía.
% NOTA CRÍTICA: el resumen de la tesis doctoral de la misma autora indica que
% la barrera más frecuente fue «prioridad» (por ejemplo, falta de tiempo), y
% que quienes ejercitaban con regularidad puntuaron más alto en disfrute,
% autoeficacia y apoyo social. La atribución del abandono «principalmente a
% la falta de guía y seguimiento personalizado» NO está respaldada por las
% fuentes revisadas; reformular el planteamiento como hipótesis de trabajo o
% citar otra fuente que sí lo sustente.
